\section{Reasoning: la piedra de fundación de la inteligencia de los modelos de lenguaje}

\subsection{Qué es Reasoning}

\begin{importantbox}{Definición central de Reasoning}
Reasoning (razonamiento) es una capacidad de «generar algo a partir de la nada»: a partir de premisas (Premise) conocidas y limitadas, se avanza a lo largo de una cadena de razonamiento para derivar conclusiones más numerosas, más ricas y más profundas.
\end{importantbox}

El silogismo clásico del razonamiento deductivo es un ejemplo típico de esto:
\begin{itemize}
    \item Premisa mayor: todos los seres humanos son mortales
    \item Premisa menor: Sócrates es un ser humano
    \item Conclusión: Sócrates también es mortal
\end{itemize}

La inteligencia real de un modelo se manifiesta precisamente a través del Reasoning. Cuando una respuesta del modelo nos deja impressed o disappointed, en esencia estamos percibiendo su capacidad de Reasoning: lo que supera la expectativa asombra, y lo que queda por debajo decepciona.

\subsection{La relación entre Reasoning y la capacidad de Agent}

\begin{knowledgebox}{Reasoning es la base de toda capacidad avanzada}
Si un modelo careciera de la capacidad de Reasoning:
\begin{itemize}
    \item No existiría \textbf{Planning}: porque Planning requiere descomponer la tarea, hacer Think Ahead y pensar varios pasos adelante
    \item No existiría \textbf{Reflection / Critique}: porque la reflexión requiere producir muchas conclusiones nuevas a partir de premisas ya existentes
    \item No existiría una conducta \textbf{Agentic} eficaz: porque un Agent necesita decidir su siguiente acción con base en el razonamiento
\end{itemize}
\end{knowledgebox}

\subsection{Dos formas clásicas de razonamiento}

\begin{enumerate}
    \item \textbf{Razonamiento deductivo (Deduction)}: razonamiento hacia adelante, de naturaleza determinista. A partir de la premisa mayor y la premisa menor se deriva una conclusión necesaria.
    \item \textbf{Razonamiento inductivo (Induction)}: razonamiento probabilístico. A partir de observaciones limitadas se induce una conclusión general; mientras más premisas existan, más confianza tiene la conclusión, pero no se garantiza que sea absolutamente correcta (como en «todos los cisnes son blancos»).
\end{enumerate}

\subsection{Reasoning Model: Long Reasoning de forma implícita}

Los Reasoning Model actuales (como DeepSeek R1 y la serie O de OpenAI) persiguen una forma de \textbf{Long Reasoning implícito}:
\begin{itemize}
    \item Problema complejo $\rightarrow$ el modelo genera espontáneamente un proceso de pensamiento largo
    \item Problema simple $\rightarrow$ genera un proceso de pensamiento corto
    \item Esta es una capacidad que el modelo aprende de manera espontánea durante la etapa de entrenamiento
\end{itemize}

\subsection{Resumen de la sección}

Reasoning es la piedra de fundación de toda capacidad avanzada de los modelos de lenguaje. Ya sea Planning, Reflection o la conducta de Agent, todas dependen de la capacidad del modelo de derivar conclusiones nuevas a partir de premisas limitadas. Los Reasoning Model modernos están reforzando esta capacidad implícita de razonamiento en cadenas largas mediante entrenamiento con RL.

